% ============================================================
%  CSCI 544 Group 36 -- Midterm Report, ACL style (acl.sty), compact layout
%  Overleaf (pdfLaTeX): upload acl.sty and acl_natbib.bst from
%  https://github.com/acl-org/acl-style-files next to this main.tex.
%  The bibliography is embedded below (filecontents).
%  NOTE: margins, title box and spacing are overridden here for space;
%  acl.sty itself is not modified.
% ============================================================
\documentclass[10pt]{article}

% ---------- Bibliography (embedded) ----------
\begin{filecontents*}[overwrite]{references.bib}
@article{magesh2025hallucination,
  title   = {Hallucination-Free? {A}ssessing the Reliability of Leading {AI} Legal Research Tools},
  author  = {Magesh, Varun and Surani, Faiz and Dahl, Matthew and Suzgun, Mirac and Manning, Christopher D. and Ho, Daniel E.},
  journal = {Journal of Empirical Legal Studies},
  volume  = {22},
  number  = {2},
  pages   = {216--242},
  year    = {2025},
  doi     = {10.1111/jels.12413}
}
@article{dahl2024legalfictions,
  title   = {Large Legal Fictions: Profiling Legal Hallucinations in Large Language Models},
  author  = {Dahl, Matthew and Magesh, Varun and Suzgun, Mirac and Ho, Daniel E.},
  journal = {Journal of Legal Analysis},
  volume  = {16},
  number  = {1},
  pages   = {64--93},
  year    = {2024},
  doi     = {10.1093/jla/laae003}
}
@inproceedings{es2024ragas,
  title     = {{RAGAs}: Automated Evaluation of Retrieval Augmented Generation},
  author    = {Es, Shahul and James, Jithin and Espinosa Anke, Luis and Schockaert, Steven},
  booktitle = {Proceedings of the 18th Conference of the European Chapter of the Association for Computational Linguistics: System Demonstrations},
  pages     = {150--158},
  year      = {2024},
  doi       = {10.18653/v1/2024.eacl-demo.16}
}
@article{laban2022summac,
  title   = {{SummaC}: Re-Visiting {NLI}-based Models for Inconsistency Detection in Summarization},
  author  = {Laban, Philippe and Schnabel, Tobias and Bennett, Paul N. and Hearst, Marti A.},
  journal = {Transactions of the Association for Computational Linguistics},
  volume  = {10},
  pages   = {163--177},
  year    = {2022},
  doi     = {10.1162/tacl_a_00453}
}
@inproceedings{barnett2024seven,
  title     = {Seven Failure Points When Engineering a Retrieval Augmented Generation System},
  author    = {Barnett, Scott and Kurniawan, Stefanus and Thudumu, Srikanth and Brannelly, Zach and Abdelrazek, Mohamed},
  booktitle = {Proceedings of the IEEE/ACM 3rd International Conference on AI Engineering -- Software Engineering for AI (CAIN)},
  year      = {2024}
}
@misc{pipitone2024legalbench,
  title  = {{LegalBench-RAG}: A Benchmark for Retrieval-Augmented Generation in the Legal Domain},
  author = {Pipitone, Nicholas and Houir Alami, Ghita},
  year   = {2024},
  note   = {arXiv:2408.10343}
}
@misc{gao2023anomalies,
  title  = {{NLP}-Based Detection of Systematic Anomalies among the Narratives of Consumer Complaints},
  author = {Gao, Peiheng and Sun, Ning and Wang, Xuefeng and Yang, Chen and Zitikis, Ri{\v{c}}ardas},
  year   = {2023},
  note   = {arXiv:2308.11138. To appear in \emph{Journal of Operational Risk}}
}
@misc{wang2026relief,
  title  = {From Complaint Narratives to Monetary Relief: A Hybrid Machine Learning Framework for {CFPB} Consumer Complaints},
  author = {Wang, Zhuoer and Zhu, Sizhen and Chen, Xiongyu},
  year   = {2026},
  note   = {arXiv:2606.22664}
}
@techreport{cfpb2026annual,
  title       = {Consumer Response Annual Report: {January} 1 -- {December} 31, 2025},
  author      = {{Consumer Financial Protection Bureau}},
  institution = {Consumer Financial Protection Bureau},
  year        = {2026},
  month       = mar,
  url         = {https://files.consumerfinance.gov/f/documents/cfpb_2025-cr-annual-report_2026-03.pdf}
}
@misc{cfpb2026process,
  title        = {Your Company's Role in the Complaint Process},
  author       = {{Consumer Financial Protection Bureau}},
  howpublished = {\url{https://www.consumerfinance.gov/compliance/consumer-complaint-program/company-process/}},
  year         = {2026},
  note         = {Accessed October 24, 2026}
}
@inproceedings{louis2022bsard,
  title     = {A Statutory Article Retrieval Dataset in {F}rench},
  author    = {Louis, Antoine and Spanakis, Gerasimos},
  booktitle = {Proceedings of the 60th Annual Meeting of the Association for Computational Linguistics (Volume 1: Long Papers)},
  pages     = {6789--6803},
  year      = {2022},
  doi       = {10.18653/v1/2022.acl-long.468}
}
@inproceedings{zheng2025reasoning,
  title     = {A Reasoning-Focused Legal Retrieval Benchmark},
  author    = {Zheng, Lucia and Guha, Neel and Arifov, Javokhir and Zhang, Sarah and Skreta, Michal and Manning, Christopher D. and Henderson, Peter and Ho, Daniel E.},
  booktitle = {Proceedings of the 4th ACM Symposium on Computer Science and Law (CSLAW)},
  year      = {2025}
}
@misc{gokhan2024rirag,
  title  = {{RIRAG}: Regulatory Information Retrieval and Answer Generation},
  author = {Gokhan, Tuba and Wang, Kexin and Gurevych, Iryna and Briscoe, Ted},
  year   = {2024},
  note   = {arXiv:2409.05677}
}
@inproceedings{ru2024ragchecker,
  title     = {{RAGChecker}: A Fine-grained Framework for Diagnosing Retrieval-Augmented Generation},
  author    = {Ru, Dongyu and Qiu, Lin and Hu, Xiangkun and Zhang, Tianhang and Shi, Peng and Chang, Shuaichen and Jiayang, Cheng and Wang, Cunxiang and Sun, Shichao and Li, Huanyu and Zhang, Zizhao and Wang, Binjie and Jiang, Jiarong and He, Tong and Wang, Zhiguo and Liu, Pengfei and Zhang, Yue and Zhang, Zheng},
  booktitle = {Advances in Neural Information Processing Systems 37 (Datasets and Benchmarks Track)},
  year      = {2024}
}
@article{liu2024lost,
  title   = {Lost in the Middle: How Language Models Use Long Contexts},
  author  = {Liu, Nelson F. and Lin, Kevin and Hewitt, John and Paranjape, Ashwin and Bevilacqua, Michele and Petroni, Fabio and Liang, Percy},
  journal = {Transactions of the Association for Computational Linguistics},
  volume  = {12},
  pages   = {157--173},
  year    = {2024},
  doi     = {10.1162/tacl_a_00638}
}
@inproceedings{gao2023alce,
  title     = {Enabling Large Language Models to Generate Text with Citations},
  author    = {Gao, Tianyu and Yen, Howard and Yu, Jiatong and Chen, Danqi},
  booktitle = {Proceedings of the 2023 Conference on Empirical Methods in Natural Language Processing},
  pages     = {6465--6488},
  year      = {2023},
  doi       = {10.18653/v1/2023.emnlp-main.398}
}
@misc{butler2026legalragbench,
  title  = {{Legal RAG Bench}: An End-to-End Benchmark for Legal {RAG}},
  author = {Butler, Abdur-Rahman and Butler, Umar},
  year   = {2026},
  note   = {arXiv:2603.01710}
}
@inproceedings{kim2025legalai,
  title     = {Where Does Legal {AI} Fail? {E}valuating {RAG} Pipelines},
  author    = {Kim, Yongjae and Lee, Wonjae},
  booktitle = {Proceedings of the 34th ACM International Conference on Information and Knowledge Management (CIKM)},
  year      = {2025},
  doi       = {10.1145/3746252.3761151}
}
@inproceedings{tang2023aggrefact,
  title     = {Understanding Factual Errors in Summarization: Errors, Summarizers, Datasets, Error Detectors},
  author    = {Tang, Liyan and Goyal, Tanya and Fabbri, Alex and Laban, Philippe and Xu, Jiacheng and Yavuz, Semih and Kryscinski, Wojciech and Rousseau, Justin and Durrett, Greg},
  booktitle = {Proceedings of the 61st Annual Meeting of the Association for Computational Linguistics (Volume 1: Long Papers)},
  pages     = {11626--11644},
  year      = {2023},
  doi       = {10.18653/v1/2023.acl-long.650}
}
@inproceedings{tang2024minicheck,
  title     = {{MiniCheck}: Efficient Fact-Checking of {LLMs} on Grounding Documents},
  author    = {Tang, Liyan and Laban, Philippe and Durrett, Greg},
  booktitle = {Proceedings of the 2024 Conference on Empirical Methods in Natural Language Processing},
  pages     = {8818--8847},
  year      = {2024},
  doi       = {10.18653/v1/2024.emnlp-main.499}
}
@inproceedings{koreeda2021contractnli,
  title     = {{ContractNLI}: A Dataset for Document-level Natural Language Inference for Contracts},
  author    = {Koreeda, Yuta and Manning, Christopher D.},
  booktitle = {Findings of the Association for Computational Linguistics: EMNLP 2021},
  pages     = {1907--1919},
  year      = {2021},
  doi       = {10.18653/v1/2021.findings-emnlp.164}
}
@article{bastani2019lda,
  title   = {Latent {D}irichlet Allocation ({LDA}) for Topic Modeling of the {CFPB} Consumer Complaints},
  author  = {Bastani, Kaveh and Namavari, Hamed and Shaffer, Jeffrey},
  journal = {Expert Systems with Applications},
  volume  = {127},
  pages   = {256--271},
  year    = {2019},
  doi     = {10.1016/j.eswa.2019.03.001}
}
@inproceedings{airani2025ensembling,
  title     = {Classification of Complaints Text Data by Ensembling Large Language Models},
  author    = {Airani, Pruthweesha and Pipada, Neha and Shah, Pratik},
  booktitle = {Proceedings of the 17th International Conference on Agents and Artificial Intelligence (ICAART)},
  volume    = {3},
  pages     = {679--686},
  year      = {2025},
  doi       = {10.5220/0013173900003890}
}
\end{filecontents*}

% ---------- ACL style ----------
\usepackage[final]{acl}   % final: no line numbers, no "Anonymous" author

\usepackage{times}
\usepackage{latexsym}
\usepackage[T1]{fontenc}
\usepackage[utf8]{inputenc}
\usepackage{microtype}
\usepackage{inconsolata}
\usepackage{graphicx}
\usepackage{booktabs}
\usepackage{multirow}
\usepackage{enumitem}
\usepackage{xcolor}
\usepackage{titlesec}
\usepackage{adjustbox}
\usepackage{xurl}

% ---------- Compact layout (overrides acl.sty defaults: 2.5cm margins, 0.6cm gutter, tall title box) ----------
\geometry{letterpaper,margin=0.55in,heightrounded=true}
\setlength\columnsep{0.22in}
\setlength\titlebox{6\baselineskip}   % one-line title + two author lines; smaller overflows into the text

\captionsetup{font=small,labelfont=bf,skip=3pt}
\setlength{\abovecaptionskip}{3pt}
\setlength{\belowcaptionskip}{0pt}
\setlength{\textfloatsep}{6pt plus 2pt minus 2pt}
\setlength{\floatsep}{6pt plus 2pt minus 2pt}
\setlength{\intextsep}{6pt plus 2pt minus 2pt}
\setlength{\dbltextfloatsep}{6pt plus 2pt minus 2pt}
\setlength{\parskip}{0pt}
\setlist{nosep,leftmargin=*}

\titleformat{\section}{\large\bfseries}{\thesection}{0.5em}{}
\titleformat{\subsection}{\normalsize\bfseries}{\thesubsection}{0.5em}{}
\titleformat{\paragraph}[runin]{\normalsize\bfseries}{}{0pt}{}[.]
\titlespacing*{\section}{0pt}{6pt plus 2pt minus 2pt}{2pt plus 1pt}
\titlespacing*{\subsection}{0pt}{4pt plus 1pt minus 1pt}{1pt plus 1pt}
\titlespacing*{\paragraph}{0pt}{2pt plus 1pt}{0.5em}

% Compact bibliography
\setlength{\bibsep}{0pt plus 0.5pt}
\renewcommand{\bibfont}{\small}

% Placeholder marker for unwritten text (delete when filled in)
\newcommand{\todo}[1]{\textcolor{gray}{[#1]}}

% Plain-text PDF title (the \adjustbox in \title is not valid PDF metadata)
\AtBeginDocument{\hypersetup{pdftitle={Regulation-Grounded Triage and Response Drafting for CFPB Consumer Complaints}}}

% ---------- Title (kept to one line) and authors ----------
\title{\adjustbox{max width=\textwidth}{Regulation-Grounded Triage and Response Drafting for CFPB Consumer Complaints}}
\author{\normalsize\mdseries Aditya Datar \quad Anirudh Joshi \quad Atharva Thorat \quad
  Gagan Shivanna \quad Ishani Bhadlikar \quad Ujjwal Puri \\
  \small\mdseries CSCI 544 Applied NLP \,\textbar\, Group 36 \,\textbar\, Midterm Report}

% ============================================================
\begin{document}
\maketitle

% ------------------------------------------------------------
\section{Introduction}

\subsection{What We Are Trying to Do}
Every company that receives a complaint through the Consumer Financial
Protection Bureau (CFPB) has 15 days to reply and up to 60 days to close the
case~\cite{cfpb2026process}. Before anyone can write that reply, someone has to
answer a hard question: which federal rule covers this complaint?

Sometimes it's easy. If you write, ``I never approved this \$240 transfer, and my
bank won't refund it,'' a trained compliance officer knows right away that
Regulation~E applies. That's the rule for electronic money transfers, and it
spells out how the bank has to investigate. Most complaints aren't that clean.
They're long. They're often angry. Parts are blacked out to protect privacy. And
they're written in everyday words, while the rules are written in legal
language.

We're building an assistant for that officer. It reads a complaint and sorts it
by product and type of problem. Then it finds the passage of federal rules that
most likely applies and drafts a short reply that cites it. Last, it checks each
sentence of the draft against that passage and flags anything the passage
doesn't back up. If nothing in our collection of rules fits, it should say so and
hand the complaint to a person. The officer always makes the final call. The
tool doesn't give legal advice.

We also want to know why drafts go wrong. Did the system find the wrong rule? Or
did it find the right one and write about it badly? So we draft every reply
twice, once from the passage our system found and once from the correct passage,
which we labeled ahead of time. Our early answer is clear. Most errors come from
finding the wrong rule, not from the writing.

\subsection{Who Cares and What Difference It Makes}\label{sec:whocares}
Compliance teams feel this first. The CFPB received about 6.6 million complaints
in 2025, up from about 3.2 million in 2024 and 1.6 million in
2023~\cite{cfpb2026annual}. The agency says part of that growth comes from credit
repair companies and from automated, often duplicate, submissions. So not every
complaint is a new case. Still, each one that reaches a company starts the
15-day clock. Someone has to read it.

Much of that reading is still done by hand. And current AI tools aren't safe to
use on their own yet. \citet{magesh2025hallucination} tested
commercial legal research tools that cite their sources, and those tools still
gave false or poorly supported answers to roughly one in six to one in three
questions. General chatbots did much worse.
\citet{dahl2024legalfictions} found they gave false answers about real
federal court cases 58\% to 88\% of the time. A citation doesn't prove an answer
is right. It doesn't tell you which sentence to doubt, either.

That's where we think our work helps. If you're the officer, a
sentence-by-sentence check lets you spend your time on the few lines that need
it. You don't have to reread everything. And if you're building a system like
ours, knowing where errors come from tells you where to put your effort: better
search, or safer writing. Our early test on 40 complaints points to search.
Given the correct passage, our automatic checker, at its default setting, found
that on average 84\% of a draft's sentences were supported and came from the
right rule. Given the passages our first search method found, only 19\% were. A
smaller sample checked by people tells the same story, about two in three
sentences against one in six. We also found no published work that links
consumer complaints to the specific regulation section that governs them. So
our labeled set of complaints may be useful to others too.

% ------------------------------------------------------------
\section{Related Work}

\subsection{Hallucination in Legal AI}
AI tools get the law wrong a lot, even when they cite sources.
\citet{magesh2025hallucination} tested three commercial legal research
tools on 202 legal questions: Lexis+ AI from LexisNexis, plus Westlaw
AI-Assisted Research and Ask Practical Law AI from Thomson Reuters. All three
look things up before they answer. All three still made the mistakes we
described in Section~\ref{sec:whocares}. \citet{dahl2024legalfictions}
saw the same problem in general chatbots answering questions about federal court
cases with no sources at all.

What matters most for us is how Magesh et al.\ define a mistake. An answer is
wrong if it's incorrect, or if it cites a source that doesn't support it or
doesn't apply. So an answer that faithfully repeats the wrong law still counts as
an error. But these studies test each tool as a whole. Magesh et al.\ couldn't
see inside the products, so they couldn't say for sure whether a mistake came
from finding the wrong source or from writing a poor answer. Our experiment is
built to answer exactly that.

\subsection{Retrieving Law from Everyday Language}
Finding the right law is hard when the question comes from an ordinary person.
BSARD~\cite{louis2022bsard} shows this. It matches 1,108 French legal questions,
based on requests from Belgian citizens, to the statutes that answer them.
Keyword search (BM25) finds 51.3\% of the relevant articles in its top 100
results. That beats neural search models used off the shelf. But it falls well
behind models trained for the task, which reach 74.8\%.
\citet{zheng2025reasoning} show why in U.S.\ law. Their questions
share very few words with the legal text that answers them, far fewer than in
common benchmarks like Natural Questions. On their housing-law questions, BM25
gets every relevant passage into its top 10 for only 18.3\% of questions, and at
least one for 40.8\%.

It looks different when the questions already sound legal. On
ObliQA~\cite{gokhan2024rirag}, where GPT-4 wrote the questions from Abu Dhabi
financial rules, BM25 comes within about 3 points of the best neural model.
LegalBench-RAG~\cite{pipitone2024legalbench}, a recent legal search benchmark,
uses contracts and privacy policies instead of laws, and it only tests the
search step. None of these benchmarks uses consumer complaints. Our guess is that
complaints sit near the everyday end, which would help explain the low BM25
scores in Section~\ref{sec:results}. We haven't measured that yet.

\subsection{Evaluating RAG and Locating Its Errors}
Our system is a form of retrieval-augmented generation (RAG). A model finds
passages first, then writes an answer from them. A common way to score these
systems is to ask whether the answer matches the passages it was given.
RAGAS~\cite{es2024ragas} is a popular example. It calls an answer faithful when
its statements follow from those passages, and that's what we mean when we call
a sentence \emph{supported}. But for legal work that isn't enough. An answer can
follow a passage perfectly and still be wrong, because the passage doesn't
apply.

Some recent work tries to split search errors from writing errors.
RAGChecker~\cite{ru2024ragchecker} scores the two steps separately, but from a
single run, by checking which retrieved passage backs each claim. Legal RAG
Bench~\cite{butler2026legalragbench} sorts each wrong answer about Australian
criminal law into a hallucination, a search error, or a reasoning error, based
on whether the right passage was found. Its conclusion is that search sets the
upper limit for many legal RAG systems. \citet{kim2025legalai} study
1,039 Korean civil complaints. They compare 12 versions of a pipeline and link
how well each one searched to how good its answers were. Handing a model the
correct passage is a common test too. Lost in the Middle~\cite{liu2024lost} and
ALCE~\cite{gao2023alce} both do it to see how well a model can do in the best
case. We run the same kind of test, holding the model fixed and changing only
the passage. Section~\ref{sec:differs} covers how ours differs. For error
analysis, we borrow the seven failure points of \citet{barnett2024seven}.
The first one, missing content, is our ``no governing section'' case.

\subsection{Checking Generated Text Against Sources}
Natural language inference (NLI) asks whether one piece of text follows from
another. It can check a draft against its source, but you have to be careful
with it. SummaC~\cite{laban2022summac} showed that NLI works better sentence by
sentence, against small chunks of the source, than against a whole document at
once. It reached 72.1 to 74.4 balanced accuracy on its benchmark, with a
separate cutoff tuned for each dataset. Our checker adapts SummaC's zero-shot
approach. It scores each sentence against blocks of the passages it cites, using
a newer model, DeBERTa-v3.

Later studies raise two warnings. These checkers do worse on text from modern
systems. On news summaries written by recent models, SummaC's zero-shot version
reaches only 63.3 balanced accuracy~\cite{tang2023aggrefact}. And a fixed cutoff
can hurt. MiniCheck~\cite{tang2024minicheck} found that SummaC's trained version
loses 9.2 points without a tuned cutoff, while MiniCheck's best small checker
comes within a point of GPT-4 for a tiny fraction of the cost. Legal text brings
its own trouble. In ContractNLI~\cite{koreeda2021contractnli}, a strong model
fails on cases where an exception written far away in the document changes the
rule. Regulations are full of phrases like ``except as provided in.'' The closest
checker built for regulations is RePASs, from RIRAG~\cite{gokhan2024rirag}. It
scores answers about Abu Dhabi rules by testing each sentence for support and
contradiction, and by counting how many required obligations the answer covers.
We took these warnings seriously. That's why we test our checker against human
labels and tune its cutoff instead of trusting the default.

\subsection{NLP on CFPB Complaints}
Most past work on CFPB complaints predicts something other than which law
applies. \citet{bastani2019lda} grouped complaint narratives into
topics and suggested the topics could help evaluate consumer protection rules.
Their topics are broad themes, though, not specific rule sections.
\citet{gao2023anomalies} predicted whether credit and prepaid card
complaints against one bank would end with relief for the consumer. It was a
step toward finding systematic patterns among complaints closed without relief,
and they used word-frequency features with classic classifiers.
\citet{wang2026relief} predicted monetary relief for 176,942 checking and
savings complaints with XGBoost and reached an AUC of 0.78. Sorting complaints by
product is common too. \citet{airani2025ensembling} report a
macro-F1 of 0.860 for DistilBERT. But that's on five product groups with a
balanced test set, so you can't compare it directly with our 11-product result.
None of these studies connects a complaint to the rule that governs it.

\subsection{How Our Work Differs}\label{sec:differs}
We found no published work that links a consumer complaint to the specific
federal regulation section that governs it. Comparing a model's answers from
retrieved and correct passages isn't new, though. ALCE~\cite{gao2023alce} does it
with the same model. What's different is how we do it.

Our correct passage is the governing section, chosen by annotators. ALCE picks
its best-case passages from the retriever's own top 100 results. We also score
every sentence twice. Is it supported by the passages it was given? And is it
supported by the section that actually applies? In the run that uses our search
results, the gap between those two scores closely matches what Magesh et al.\
call misgrounded. And we do all of this sentence by sentence on federal
regulations, with a checker we've tested against human labels.

% ------------------------------------------------------------
\section{Data and Annotation}

\subsection{Consumer Complaints}
\todo{FOIA narratives archive, cleaning funnel, splits.}

\subsection{Regulation Corpus}
\todo{Regulations E, F, V, Z and the FCRA statute; structure-aware chunking.}

\subsection{Gold Annotation}
\todo{198 complaints, LLM annotators with human adjudication, agreement.}

% ------------------------------------------------------------
\section{Method}

\subsection{System Overview}
\todo{Figure 1: the four-stage pipeline and human handoff.}

\paragraph{Classification}
\todo{DistilBERT, product and issue.}

\paragraph{Retrieval}
\todo{BM25 baseline; dense and hybrid planned.}

\paragraph{Response drafting}
\todo{Nemotron with per-sentence passage tags.}

\paragraph{Sentence checking}
\todo{DeBERTa-v3 NLI following SummaC.}

\subsection{Attribution Experiment}
\todo{Retrieved vs.\ correct passage; ``supported'' vs.\ ``supported and correct''.}

\subsection{What Is New}
\todo{Three contributions.}

% ------------------------------------------------------------
\section{Experimental Plan}

\subsection{Hypotheses}
\todo{H1--H5 (R6).}

\subsection{Evaluation}
\todo{Metrics per stage and end to end (R7).}

% ------------------------------------------------------------
\section{Preliminary Results}\label{sec:results}

\subsection{Classification}
\todo{Product and issue macro-F1.}

\subsection{Retrieval}
\todo{BM25 hit@k and MRR; abstention.}

\subsection{Checker Validation}
\todo{AUROC, cutoff sensitivity, agreement with human labels.}

\subsection{Attribution Pilot}
\todo{Retrieved vs.\ correct passage gap, human-label estimate.}

\subsection{Limitations So Far}
\todo{Pilot size, silver labels, metric definition.}

% ------------------------------------------------------------
\section{Milestones}

\subsection{Progress Against the Proposal}
\todo{M1 and M2 status, and changes from the proposal (R9).}

\subsection{Updated Milestones}
\todo{Plan for the rest of the semester (R10).}

% ------------------------------------------------------------
\bibliography{references}

\end{document}
